\documentclass[11pt,letterpaper]{article}
\usepackage[utf8]{inputenc}
\usepackage[T1]{fontenc}
\usepackage{geometry}
\geometry{margin=1in}
\usepackage{hyperref}

\pagestyle{empty}

\begin{document}

\today

\vspace{1cm}

Dear Editor,

\vspace{0.5cm}

We submit for your consideration the manuscript entitled ``\textbf{Context-Dependent Expression Persistence: AR(2) Eigenvalue $|\lambda|$ Is Independent of Intrinsic mRNA Half-Life}'' for publication in \textit{Genome Biology}.

\vspace{0.3cm}

\textbf{Summary.} We demonstrate that the eigenvalue modulus of a second-order autoregressive model---a measure of temporal persistence in gene expression---is independent of intrinsic mRNA half-life ($\rho = 0.006$, $n = 5{,}945$ genes). This near-zero correlation, validated across four non-circadian datasets spanning three species, establishes that temporal persistence (how strongly a gene's past expression determines its future) is a genuinely distinct biological property from transcript stability.

\vspace{0.3cm}

\textbf{Significance.} A critical challenge in interpreting any temporal metric applied to gene expression is distinguishing genuine regulatory dynamics from trivial biochemical properties. We show that the AR(2) eigenvalue modulus captures context-dependent regulatory persistence---arising from transcription factor dynamics, feedback loops, and signaling cascades---rather than the intrinsic decay rate of mRNA molecules. The IFIT1 case study (31-minute mRNA half-life, $|\lambda| = 0.72$) vividly illustrates this dissociation: a gene with one of the shortest-lived transcripts in the genome shows high temporal persistence due to sustained interferon-driven retranscription. This work is a companion to our Paper A (submitted to \textit{PLOS Computational Biology}), which establishes the eigenvalue modulus as a measure of circadian hierarchy.

\vspace{0.3cm}

\textbf{Reproducibility.} All datasets are publicly available from NCBI GEO and published databases. The PAR(2) Discovery Engine web application (\url{https://par2-discovery-engine.replit.app}) provides interactive access to all results and allows users to upload their own data. A standalone reproducibility package is included in the Supplementary Materials.

\vspace{0.3cm}

\textbf{Limitations acknowledged.} We are transparent about the method's boundaries: mRNA half-life data comes from embryonic stem cells (not liver), static actinomycin D measurements may miss dynamically regulated stability, and the AR(2) model assumes linearity and stationarity. Despite these caveats, the near-zero correlation is robust across all datasets and analysis approaches tested.

\vspace{0.3cm}

This manuscript has not been submitted elsewhere and all work is original.

\vspace{0.5cm}

Sincerely,

\vspace{0.5cm}

Michael Whiteside \\
Independent Researcher, United Kingdom \\
mickwh@msn.com \\
ORCID: 0009-0000-0643-5791

\end{document}